% !TEX program = xelatex
\documentclass[12pt]{article}
\usepackage[UTF8]{ctex}
\usepackage{booktabs}
\usepackage{amsmath}
\usepackage{amssymb}
\usepackage{geometry}
\usepackage{caption}
\usepackage{threeparttable}
\usepackage{float}
\usepackage{setspace}
\usepackage{pdflscape}
\usepackage{multirow}
\usepackage{array}
\usepackage{enumitem}

\geometry{left=1.1in, right=1.1in, top=1in, bottom=1in}
\setstretch{1.2}
\captionsetup{labelfont=bf, font=small}

\newcommand{\sym}[1]{\ensuremath{^{\text{#1}}}}
\newcommand{\se}[1]{{\small(#1)}}
\newcommand{\abs}{\textemdash}
\newcommand{\mis}{}

% 自动换行列类型
\newcolumntype{L}[1]{>{\raggedright\arraybackslash}p{#1}}

\begin{document}

% ============================================================
%  封面
% ============================================================
\begin{center}
  {\large\bfseries 外包采购价格的决定因素：回归结果汇总}\\[6pt]
  {\normalsize 2026-05-29}
\end{center}

\medskip
\noindent\textbf{数据与样本：}
2017年全国增值税发票数据。样本限定条件：
\begin{enumerate}[nosep, leftmargin=2em]
  \item \textbf{外包产品}：同一企业在同一年内对同一产品编码（9位）既有采购记录又有销售记录，外包额 $=\min(\text{采购额},\,\text{销售额})>0$。
  \item \textbf{非中介企业}：剔除外包比例 $>90\%$ 的纯转卖商。
\end{enumerate}
主回归样本（表二至表四）共 $N\approx46{,}100$ 个企业$\times$产品观测。表一额外包含资产变量，受汇算清缴覆盖率（约48\%）限制，$N\approx19{,}200$。

\bigskip

% ============================================================
%  TABLE 1
% ============================================================
\section*{表一：基准回归（逐步加入固定效应）}

\subsection*{回归方程}

\begin{align}
  \ln p_{fjct}
    &= \gamma_1\ln\mathrm{Output}_{ft}
     + \gamma_2\ln K_{ft}
     + \gamma_3\,n\mathrm{Prod}_{ft} \notag\\
    &\quad + \lambda_1\ln n^{B}_{jct}
           + \lambda_2\ln n^{S}_{jct}
           + \lambda_3\ln q^{m}_{jct}
           + \lambda_4\ln\bar{p}^{m}_{jct} \notag\\
    &\quad + \delta_{f}^{[k]} + \delta_{j}^{[k]} + \delta_{c}^{[k]}
           + \varepsilon_{fjct}
  \label{eq:baseline}
\end{align}

\noindent 上标 $[k]$ 表示各列逐步加入固定效应（见表格底部）。企业特征变量（$\gamma_1,\gamma_2,\gamma_3$）在单年截面中被企业固定效应完全吸收，仅在 OLS 规格（列1--2）可识别。

\medskip
\noindent\textbf{变量定义}

\smallskip
{\small
\begin{tabular}{@{} L{3.5cm} L{10.5cm} @{}}
\toprule
符号 & 含义 \\
\midrule
$p_{fjct}$ &
  \textbf{因变量}：企业 $f$ 在城市 $c$、年份 $t$ 外包采购产品 $j$ 的
  加权平均单价（元/单位），取对数。$f$=企业，$j$=产品（9位商品编码），
  $c$=地级市（买方所在城市），$t$=年份（本文2017年截面）。\\[4pt]
$\ln\mathrm{Output}_{ft}$ &
  企业 $f$ 年总增值税销售额对数，作为企业规模代理变量。\\[4pt]
$\ln K_{ft}$ &
  企业 $f$ 资产总额对数，覆盖率约48\%，其余观测缺失。\\[4pt]
$n\mathrm{Prod}_{ft}$ &
  企业 $f$ 年度生产产品种数（净生产额$>0$的产品数）。\\[4pt]
$\ln n^{B}_{jct}$ &
  城市 $c$ 中购买产品 $j$ 的企业数对数，衡量买方市场厚度。\\[4pt]
$\ln n^{S}_{jct}$ &
  城市 $c$ 中销售产品 $j$ 的企业数对数，衡量供给侧竞争程度。\\[4pt]
$\ln q^{m}_{jct}$ &
  城市 $c$ 中产品 $j$ 的市场总采购量对数。\\[4pt]
$\ln\bar{p}^{m}_{jct}$ &
  城市 $c$ 中产品 $j$ 的金额加权市场均价对数，作为外部价格锚点。\\[4pt]
$\delta_{f}^{[k]},\,\delta_{j}^{[k]},\,\delta_{c}^{[k]}$ &
  企业、产品、城市固定效应，逐步加入。
  因企业与城市一一对应，City FE 被 Firm FE 完全吸收（列4$=$列5）。\\
\bottomrule
\end{tabular}
}

\vspace{8pt}

\begin{table}[H]
\centering\small
\begin{threeparttable}
\caption{基准回归：逐步加入固定效应}
\begin{tabular}{l *{5}{c}}
\toprule
 & (1) & (2) & (3) & (4) & (5) \\
 & OLS-bare & OLS-full & {+Firm FE} & {+Firm+Prod} & {+Firm+Prod+City} \\
\midrule
$\ln\mathrm{Output}_{ft}$
  & $0.112\sym{***}$ & $0.087\sym{***}$ & \abs & \abs & \abs \\
  & \se{0.029} & \se{0.018} & & & \\[3pt]
$\ln K_{ft}$
  & $-0.014$ & $-0.016$ & \abs & \abs & \abs \\
  & \se{0.027} & \se{0.018} & & & \\[3pt]
$n\mathrm{Prod}_{ft}$
  & $-0.001\sym{**}$ & $-0.001\sym{***}$ & \abs & \abs & \abs \\
  & \se{0.000} & \se{0.000} & & & \\[3pt]
$\ln n^{B}_{jct}$
  & \mis & $0.236\sym{***}$ & $0.060$ & $0.019$ & $0.019$ \\
  & & \se{0.057} & \se{0.043} & \se{0.051} & \se{0.051} \\[3pt]
$\ln n^{S}_{jct}$
  & \mis & $-0.040$ & $0.000$ & $-0.134\sym{***}$ & $-0.134\sym{***}$ \\
  & & \se{0.072} & \se{0.057} & \se{0.050} & \se{0.050} \\[3pt]
$\ln q^{m}_{jct}$
  & \mis & $-0.012$ & $0.015$ & $-0.008$ & $-0.008$ \\
  & & \se{0.020} & \se{0.020} & \se{0.025} & \se{0.025} \\[3pt]
$\ln\bar{p}^{m}_{jct}$
  & \mis & $0.771\sym{***}$ & $0.663\sym{***}$ & $0.191\sym{***}$ & $0.191\sym{***}$ \\
  & & \se{0.028} & \se{0.027} & \se{0.044} & \se{0.044} \\[3pt]
Constant
  & $2.532\sym{***}$ & $-0.622$ & $1.519\sym{***}$ & $4.350\sym{***}$ & $4.350\sym{***}$ \\
  & \se{0.640} & \se{0.561} & \se{0.300} & \se{0.454} & \se{0.456} \\
\midrule
Firm FE    & No  & No  & Yes & Yes & Yes \\
Product FE & No  & No  & No  & Yes & Yes \\
City FE    & No  & No  & No  & No  & Yes \\
$N$        & 19,735 & 19,701 & 19,549 & 19,195 & 19,195 \\
Adj.\ $R^{2}$ & 0.018 & 0.349 & 0.482 & 0.624 & 0.621 \\
\bottomrule
\end{tabular}
\begin{tablenotes}\small
\item \textit{注：}``\textemdash'' 表示被固定效应吸收，系数不可识别。所有列因含 $\ln K_{ft}$ 而受限于汇算清缴覆盖率（$N\approx19{,}200$）。$\ln K_{ft}$ 在所有列均不显著。$\ln n^{S}_{jct}$ 在加入产品固定效应后显著为负（$-0.134^{***}$），表明控制产品质量异质性后供给侧竞争显著压低采购价格。
\end{tablenotes}
\end{threeparttable}
\end{table}

\clearpage

% ============================================================
%  TABLE 2
% ============================================================
\section*{表二：需求侧 vs 供给侧分解}

\subsection*{回归方程}

\begin{align}
  \ln p_{fjct}
    &= \underbrace{\lambda_1\ln n^{B}_{jct}
                 + \lambda_3\ln q^{m}_{jct}}_{\text{需求侧}}
     + \underbrace{\lambda_2\ln n^{S}_{jct}}_{\text{供给侧}}
     + \underbrace{\lambda_4\ln\bar{p}^{m}_{jct}}_{\text{市场均价}} \notag\\
    &\quad + \delta_{f} + \delta_{j} + \delta_{c} + \varepsilon_{fjct}
  \label{eq:demand_supply}
\end{align}

\noindent 全程保留 Firm + Product + City 固定效应，各列逐步加入需求侧、供给侧、市场均价变量。

\medskip
\noindent\textbf{变量定义}

\smallskip
{\small
\begin{tabular}{@{} L{3.5cm} L{10.5cm} @{}}
\toprule
符号 & 含义 \\
\midrule
$\ln n^{B}_{jct}$，$\ln q^{m}_{jct}$ &
  \textbf{需求侧变量}：买方数与市场总采购量，刻画买方侧竞争结构。
  更多买方或更大市场规模预期对应更高价格（卖方市场议价力更强）。\\[4pt]
$\ln n^{S}_{jct}$ &
  \textbf{供给侧变量}：卖方数，刻画供给竞争程度。
  更多卖方预期对应更低价格（买方议价力更强）。\\[4pt]
$\ln\bar{p}^{m}_{jct}$ &
  \textbf{市场均价}：同城市同产品的外部价格锚点，综合反映供需均衡信息。
  加入后预期其余竞争度变量系数收缩至不显著。\\[4pt]
其余符号 & 同表一。\\
\bottomrule
\end{tabular}
}

\vspace{8pt}

\begin{table}[H]
\centering\small
\begin{threeparttable}
\caption{需求侧 vs 供给侧分解}
\begin{tabular}{l *{4}{c}}
\toprule
 & (1) & (2) & (3) & (4) \\
 & Demand & Supply & Both & Both$+\bar{p}^{m}$ \\
\midrule
$\ln n^{B}_{jct}$
  & $0.104\sym{***}$ & \mis & $0.107\sym{***}$ & $0.022$ \\
  & \se{0.032} & & \se{0.034} & \se{0.036} \\[3pt]
$\ln q^{m}_{jct}$
  & $-0.106\sym{***}$ & \mis & $-0.106\sym{***}$ & $-0.003$ \\
  & \se{0.013} & & \se{0.013} & \se{0.015} \\[3pt]
$\ln n^{S}_{jct}$
  & \mis & $-0.065\sym{**}$ & $-0.006$ & $-0.051$ \\
  & & \se{0.031} & \se{0.033} & \se{0.034} \\[3pt]
$\ln\bar{p}^{m}_{jct}$
  & \mis & \mis & \mis & $0.158\sym{***}$ \\
  & & & & \se{0.024} \\[3pt]
Constant
  & $5.141\sym{***}$ & $4.457\sym{***}$ & $5.151\sym{***}$ & $3.792\sym{***}$ \\
  & \se{0.235} & \se{0.163} & \se{0.244} & \se{0.268} \\
\midrule
Firm FE    & Yes & Yes & Yes & Yes \\
Product FE & Yes & Yes & Yes & Yes \\
City FE    & Yes & Yes & Yes & Yes \\
$N$        & 46,163 & 46,106 & 46,106 & 46,106 \\
Adj.\ $R^{2}$ & 0.601 & 0.599 & 0.601 & 0.603 \\
\bottomrule
\end{tabular}
\begin{tablenotes}\small
\item \textit{注：}列(1)与列(2)样本差异（46,163 vs 46,106）源于 $\ln n^{S}_{jct}$ 有少量缺失。不含市场均价时，需求侧（$\ln n^{B}=0.107^{***}$，$\ln q^{m}=-0.106^{***}$）与供给侧（$\ln n^{S}=-0.065^{**}$）均显著且方向符合预期。加入 $\ln\bar{p}^{m}_{jct}$（列4）后所有变量失显，表明市场均价已包含竞争信息的主要成分。
\end{tablenotes}
\end{threeparttable}
\end{table}

\clearpage

% ============================================================
%  TABLE 3
% ============================================================
\section*{表三：产品相似度（$S_{mj}$，$C_{mj}$）}

\subsection*{回归方程}

\begin{align}
  \ln p_{fjct}
    &= \lambda_1\ln n^{B}_{jct}
     + \lambda_2\ln n^{S}_{jct}
     + \lambda_3\ln q^{m}_{jct}
     + \lambda_4\ln\bar{p}^{m}_{jct} \notag\\
    &\quad + \beta_1 S_{mj} + \beta_2 C_{mj} \notag\\
    &\quad + \delta_{f} + \delta_{j} + \delta_{c} + \varepsilon_{fjct}
  \label{eq:similarity}
\end{align}

\noindent 列(1)为基准规格（无相似度），列(2)--(4)逐步加入 $S_{mj}$、$C_{mj}$。

\medskip
\noindent\textbf{变量定义}

\smallskip
{\small
\begin{tabular}{@{} L{3.5cm} L{10.5cm} @{}}
\toprule
符号 & 含义 \\
\midrule
$S_{mj}$ &
  \textbf{投入结构相似度}：企业核心产品 $m$ 与外包产品 $j$
  在投入品结构上的余弦相似度。$S_{mj}$ 较高表示企业对 $j$ 的
  供应链更熟悉，预期议价能力更强，采购价格更低（系数预期为负）。\\[4pt]
$C_{mj}$ &
  \textbf{产出结构相似度}：企业核心产品 $m$ 与外包产品 $j$
  在下游客户/分销渠道上的重叠度。$C_{mj}$ 较高表示产品 $j$
  与主业关联更紧密，采购价格影响方向待检验。\\[4pt]
$m$ &
  \textbf{企业核心产品}：企业年度净生产额（$=$销售额$-$同产品采购额）
  最大的产品。$S_{mj}$ 和 $C_{mj}$ 以 $m$ 为基准预先计算。\\[4pt]
$\beta_1,\,\beta_2$ &
  相似度系数。$S_{mj}$、$C_{mj}$ 在 firm$\times$product 层面变化，
  不被企业固定效应或产品固定效应单独吸收，FE 规格下可识别。\\[4pt]
其余符号 & 同表一。\\
\bottomrule
\end{tabular}
}

\vspace{8pt}

\begin{table}[H]
\centering\small
\begin{threeparttable}
\caption{产品相似度：$S_{mj}$（投入）与 $C_{mj}$（产出）}
\begin{tabular}{l *{4}{c}}
\toprule
 & (1) & (2) & (3) & (4) \\
 & Baseline & $+S_{mj}$ & $+C_{mj}$ & Both \\
\midrule
$\ln n^{B}_{jct}$
  & $0.022$ & $0.022$ & $0.022$ & $0.022$ \\
  & \se{0.036} & \se{0.036} & \se{0.036} & \se{0.036} \\[3pt]
$\ln n^{S}_{jct}$
  & $-0.051$ & $-0.051$ & $-0.051$ & $-0.051$ \\
  & \se{0.034} & \se{0.034} & \se{0.034} & \se{0.034} \\[3pt]
$\ln q^{m}_{jct}$
  & $-0.003$ & $-0.003$ & $-0.003$ & $-0.003$ \\
  & \se{0.015} & \se{0.015} & \se{0.015} & \se{0.015} \\[3pt]
$\ln\bar{p}^{m}_{jct}$
  & $0.158\sym{***}$ & $0.158\sym{***}$ & $0.158\sym{***}$ & $0.158\sym{***}$ \\
  & \se{0.024} & \se{0.024} & \se{0.024} & \se{0.024} \\[3pt]
$S_{mj}$
  & \mis & $0.002$ & \mis & $-0.006$ \\
  & & \se{0.075} & & \se{0.090} \\[3pt]
$C_{mj}$
  & \mis & \mis & $0.007$ & $0.013$ \\
  & & & \se{0.090} & \se{0.108} \\[3pt]
Constant
  & $3.792\sym{***}$ & $3.793\sym{***}$ & $3.794\sym{***}$ & $3.794\sym{***}$ \\
  & \se{0.268} & \se{0.268} & \se{0.267} & \se{0.268} \\
\midrule
Firm FE    & Yes & Yes & Yes & Yes \\
Product FE & Yes & Yes & Yes & Yes \\
City FE    & Yes & Yes & Yes & Yes \\
$N$        & 46,106 & 46,094 & 46,094 & 46,094 \\
Adj.\ $R^{2}$ & 0.603 & 0.603 & 0.603 & 0.603 \\
\bottomrule
\end{tabular}
\begin{tablenotes}\small
\item \textit{注：}加入相似度后样本仅减少12个观测值（覆盖率$>99.9\%$），表明外包产品在 \texttt{full\_data.dta} 中几乎完整覆盖。$S_{mj}$ 和 $C_{mj}$ 系数接近零且不显著，$R^{2}$ 无变化，表明在控制市场均价后产品相似度与外包采购价格无显著关联。
\end{tablenotes}
\end{threeparttable}
\end{table}

\clearpage

% ============================================================
%  TABLE 4  —  方程与变量说明（portrait），回归表（landscape）
% ============================================================
\section*{表四：企业规模 $\times$ 市场条件交互}

\subsection*{回归方程}

\begin{align}
  \ln p_{fjct}
    &= \lambda_1\ln n^{B}_{jct}
     + \lambda_2\ln n^{S}_{jct}
     + \lambda_3\ln q^{m}_{jct}
     + \lambda_4\ln\bar{p}^{m}_{jct} \notag\\
    &\quad
     + \mu_1\bigl(\ln\mathrm{Output}_{ft}\times\ln n^{B}_{jct}\bigr)
     + \mu_2\bigl(\ln\mathrm{Output}_{ft}\times\ln n^{S}_{jct}\bigr) \notag\\
    &\quad
     + \mu_3\bigl(\ln\mathrm{Output}_{ft}\times\ln\bar{p}^{m}_{jct}\bigr)
     + \mu_4\bigl(\ln\mathrm{Output}_{ft}\times\ln q^{m}_{jct}\bigr) \notag\\
    &\quad
     + \delta_{f} + \delta_{j} + \delta_{c} + \varepsilon_{fjct}
  \label{eq:interactions}
\end{align}

\noindent 各列逐步加入单个交互项，列(5)全部加入。$\ln\mathrm{Output}_{ft}$ 主效应被 Firm FE 吸收（系数显示为``\textemdash''），但交互项在 firm$\times$product 层面变化，FE 规格下可识别。

\medskip
\noindent\textbf{变量定义}

\smallskip
{\small
\begin{tabular}{@{} L{4.2cm} L{9.8cm} @{}}
\toprule
符号 & 含义 \\
\midrule
$\ln\mathrm{Output}_{ft}$ &
  企业规模代理：企业年总增值税销售额对数（同表一），
  主效应被 Firm FE 吸收，仅交互项可识别。\\[4pt]
$\ln\mathrm{Output}\times\ln n^{B}$ &
  规模$\times$买方厚度：检验大企业在买方更多的市场中议价优势是否更强。
  预期系数为负（大企业在厚买方市场中拿到更低价格）。\\[4pt]
$\ln\mathrm{Output}\times\ln n^{S}$ &
  规模$\times$卖方竞争：检验大企业在供给竞争激烈时议价优势是否更强。
  预期系数为负。\\[4pt]
$\ln\mathrm{Output}\times\ln\bar{p}^{m}$ &
  规模$\times$市场均价：检验大企业在价格水平较高的市场中能否相对节约采购成本。\\[4pt]
$\ln\mathrm{Output}\times\ln q^{m}$ &
  规模$\times$市场规模：检验大企业是否更能利用市场总量优势压低价格。
  预期系数为负。\\[4pt]
$\mu_1,\,\mu_2,\,\mu_3,\,\mu_4$ &
  交互项系数，捕捉企业规模对市场条件定价效应的异质性调节作用。\\[4pt]
其余符号 & 同表一。\\
\bottomrule
\end{tabular}
}

% 回归表置于横向页
\begin{landscape}
\vspace*{\fill}
\begin{table}[H]
\centering\small
\begin{threeparttable}
\caption{企业规模 $\times$ 市场条件交互}
\begin{tabular}{l *{5}{c}}
\toprule
 & (1) & (2) & (3) & (4) & (5) \\
 & Baseline
 & {$+$Size$\times n^{B}$}
 & {$+$Size$\times n^{S}$}
 & {$+$Size$\times\bar{p}^{m}$}
 & All \\
\midrule
$\ln\mathrm{Output}_{ft}$
  & \abs & \abs & \abs & \abs & \abs \\
$n\mathrm{Prod}_{ft}$
  & \abs & \abs & \abs & \abs & \abs \\[3pt]
$\ln n^{B}_{jct}$
  & $0.022$ & $0.413\sym{***}$ & $0.026$ & $0.029$ & $0.018$ \\
  & \se{0.036} & \se{0.101} & \se{0.036} & \se{0.036} & \se{0.137} \\[3pt]
$\ln n^{S}_{jct}$
  & $-0.051$ & $-0.057\sym{*}$ & $0.415\sym{***}$ & $-0.053$ & $0.132$ \\
  & \se{0.034} & \se{0.034} & \se{0.115} & \se{0.033} & \se{0.145} \\[3pt]
$\ln q^{m}_{jct}$
  & $-0.003$ & $-0.004$ & $-0.005$ & $-0.007$ & $0.131\sym{*}$ \\
  & \se{0.015} & \se{0.015} & \se{0.015} & \se{0.015} & \se{0.070} \\[3pt]
$\ln\bar{p}^{m}_{jct}$
  & $0.158\sym{***}$ & $0.157\sym{***}$ & $0.156\sym{***}$ & $-0.159\sym{**}$ & $-0.015$ \\
  & \se{0.024} & \se{0.024} & \se{0.024} & \se{0.080} & \se{0.073} \\[3pt]
$\ln\mathrm{Output}\times\ln n^{B}$
  & \mis & $-0.019\sym{***}$ & \mis & \mis & $0.001$ \\
  & & \se{0.005} & & & \se{0.006} \\[3pt]
$\ln\mathrm{Output}\times\ln n^{S}$
  & \mis & \mis & $-0.024\sym{***}$ & \mis & $-0.009$ \\
  & & & \se{0.006} & & \se{0.007} \\[3pt]
$\ln\mathrm{Output}\times\ln\bar{p}^{m}$
  & \mis & \mis & \mis & $0.016\sym{***}$ & $0.009\sym{**}$ \\
  & & & & \se{0.004} & \se{0.004} \\[3pt]
$\ln\mathrm{Output}\times\ln q^{m}$
  & \mis & \mis & \mis & \mis & $-0.007\sym{**}$ \\
  & & & & & \se{0.004} \\[3pt]
Constant
  & $3.792\sym{***}$ & $3.767\sym{***}$ & $3.785\sym{***}$ & $3.836\sym{***}$ & $3.811\sym{***}$ \\
  & \se{0.268} & \se{0.268} & \se{0.265} & \se{0.266} & \se{0.264} \\
\midrule
Firm FE    & Yes & Yes & Yes & Yes & Yes \\
Product FE & Yes & Yes & Yes & Yes & Yes \\
City FE    & Yes & Yes & Yes & Yes & Yes \\
$N$        & 46,106 & 46,106 & 46,106 & 46,106 & 46,106 \\
Adj.\ $R^{2}$ & 0.603 & 0.603 & 0.603 & 0.604 & 0.604 \\
\bottomrule
\end{tabular}
\begin{tablenotes}\small
\item \textit{注：}``\textemdash'' 表示被 Firm FE 吸收。单项交互（列2--4）均显著：大企业在买方/卖方更多的市场中采购价格更低（$-0.019^{***}$，$-0.024^{***}$），在价格水平高的市场中支付相对更高价格（$+0.016^{***}$）。全部加入（列5）后 $\mathrm{Size}\times\bar{p}^{m}$（$0.009^{**}$）和 $\mathrm{Size}\times q^{m}$（$-0.007^{**}$）同时显著，大企业在价格高、规模大的市场中的差异化议价效应稳健。
\end{tablenotes}
\end{threeparttable}
\end{table}
\vspace*{\fill}
\end{landscape}

\end{document}
